\section{Tabela prawdy}
\label{sec:tabela}

Tabela prawdy dla wyrażenia~\eqref{eq:sop} zawiera $2^3 = 8$ wierszy
(trzy zmienne wejściowe: $Q$, $P$, $T$).

\begin{table}[H]
\centering
\caption{Tabela prawdy układu.}
\label{tab:truth}
\begin{tabular}{ccc|c|l}
\toprule
$Q$ & $P$ & $T$ & $Q_{\text{next}}$ & Interpretacja \\
\midrule
0 & 0 & 0 & 0 & ciemno, brak przycisku, brak przekroczenia czasu \\
0 & 0 & 1 & 0 & ciemno, bez zmian \\
0 & 1 & 0 & 1 & \textbf{włączenie} światła przyciskiem \\
0 & 1 & 1 & 0 & przycisk wciśnięty, ale timeout gasi światło \\
1 & 0 & 0 & 1 & świeci, brak przycisku, brak timeoutu \\
1 & 0 & 1 & 0 & \textbf{automatyczne wyłączenie} po 2 min bezruchu \\
1 & 1 & 0 & 0 & \textbf{wyłączenie} światła przyciskiem \\
1 & 1 & 1 & 0 & timeout wygrywa nad przyciskiem \\
\bottomrule
\end{tabular}
\end{table}

\paragraph{Uwaga projektowa.} W wierszach, w których $T=1$, wyjście
jest zerowane niezależnie od stanu $P$ i $Q$ – jest to zamierzone:
sygnał timeoutu ma najwyższy priorytet (asynchroniczne kasowanie
przerzutnika).